\section{Methodology}

\subsection{System Architecture}

Our system instantiates a small panel of AI-simulated domain experts to evaluate puzzle hunt puzzles using a structured scoring rubric. The pipeline proceeds in five stages: (1) a \emph{profile library} of persistent expert reviewer profiles is maintained as version-controlled markdown artifacts; (2) for each puzzle under evaluation, a \emph{panel assembly} step selects three reviewers from the library with complementary evaluative lenses (e.g., structural rigor paired with accessibility and narrative); (3) each selected profile is loaded as full system context and the model is prompted to generate a written \emph{review} from that expert's perspective; (4) the three reviews are \emph{synthesized} to surface consensus findings, unresolved disagreements, and actionable revisions; (5) a \emph{score} on six or seven dimensions is extracted (see Section~\ref{sec:rubric}), producing a per-puzzle quality profile (see Figure~1 for the full pipeline diagram). All experiments reported here use Claude Sonnet~4.6 as the underlying model, with temperature 1.0 and no output truncation. Profiles are persistent artifacts --- they are authored once, refined over time, and loaded unchanged for each review session. No profile content is regenerated or conditioned on the puzzle being evaluated.

\subsection{Profile Design}

Each reviewer profile is a markdown document of 80--120 lines organized into six sections: \textbf{Identity} (name, affiliation, and credentials as a puzzle or game designer); \textbf{Design Philosophy} (3--4 paragraphs articulating the reviewer's core convictions in first-person or closely attributed voice); \textbf{Review Lens} (a structured list of 10--14 specific evaluative questions the reviewer applies to every puzzle); \textbf{Key Works} (titles and roles that ground the persona in documented output); \textbf{Puzzle Hunt Credentials} (competitive and constructive history relevant to the hunt domain); and \textbf{Key Sources} (links and publications anchoring the profile to verifiable public record).

The \emph{review lens} section is the profile's primary functional contribution: it operationalizes the reviewer's philosophy as a checklist of domain-specific probes that the model executes during evaluation. Table~1 summarizes the structural differences between the two profile versions used in our experiments.

\begin{table}[h]
\centering
\caption{Profile version comparison. Named frameworks are evaluative concepts introduced in the review lens with specific names (e.g., \emph{load-bearing test}, \emph{a-ha economy}, \emph{world-model consistency}).}
\begin{tabular}{lcccc}
\hline
Version & Profiles & Avg.\ lines & Named frameworks & Characteristic concerns \\
\hline
v1 & 12 & $\sim$50 & No & No \\
v2 & 29 & $\sim$95 & Yes & Yes \\
\hline
\end{tabular}
\label{tab:profiles}
\end{table}

The key design decision distinguishing v1 from v2 is the introduction of \emph{named evaluative frameworks} --- domain-specific analytical concepts that the reviewer applies by name, enabling the model to produce critique that is transferable and citable rather than generic. To make this concrete, consider the same reviewer (Thomas Snyder) across both versions. In v1, the analogous profile entry reads: ``Checks that each element is necessary and that the puzzle has a clear intended solve path.'' In v2, the corresponding review lens asks: ``Is each element load-bearing? He strips individual constraints, given values, or visual elements and tests whether the puzzle still forces a unique solution. Anything that can be removed without consequence was not design --- it was accumulation.'' The v2 formulation names a test (\emph{load-bearing test}), specifies the operation (remove and check uniqueness), and characterizes failure in terms the reviewer would use in actual feedback. Similar named frameworks introduced in v2 include the \emph{a-ha economy} (does the puzzle produce one satisfying moment of insight, or squander it through redundancy), \emph{world-model consistency} (does the puzzle's logic hold within the fictional or factual universe it invokes), \emph{social fabric} (does the puzzle work in a team-solving context or require solitary reconstruction), and \emph{perceptual-shift} (does the solver experience a genuine reframing of the problem, or merely accumulate information). The presence of these frameworks in review output, and their absence in v1 output, is the primary signal we measure in Experiment~0.

\subsection{Scoring Rubric}
\label{sec:rubric}

Puzzles are scored on seven dimensions by each of the three panel reviewers. The rubric design is itself an empirical contribution of this paper: we begin with a six-dimension equal-weight baseline (used in Experiments~0--2 to control for rubric effects while testing profile depth), then validate a seven-dimension weighted configuration in Section~4.4. Table~2 provides full definitions for all seven dimensions.

\textbf{Clarity} (×1): whether a solver can understand what the puzzle is asking without external help. Hunt puzzles intentionally withhold genre labels; the mechanism must do all signposting.

\textbf{Solvability} (×1): whether a solver with the target domain knowledge can reliably reach the correct answer. Distinct from difficulty: a puzzle can be challenging and fully solvable, or easy and broken.

\textbf{Elegance} (×2 in validated rubric): whether the puzzle's mechanism feels clean, satisfying, and non-arbitrary. Elegance signals intentional design; we distinguish elegant puzzles from merely correct ones. Double-weighted in the validated rubric because it is empirically the strongest tier discriminator across 15 puzzle types.

\textbf{Reading Reward} (×2 in validated rubric): whether domain knowledge is load-bearing. A puzzle solvable by general reasoning without domain knowledge has failed its primary design criterion. Double-weighted alongside Elegance; both are proxies for the Riven Standard (see below).

\textbf{Fun} (×1): whether the experience is enjoyable and produces an identifiable \emph{aha} moment. The primary phenomenological output of puzzle design.

\textbf{Confirmation} (×1): whether the solver knows when they have reached the correct answer without being told. Weak confirmation produces anxiety and hedged submissions.

\textbf{Riven Standard} (×1, validated rubric only): the field's highest-order quality criterion, named after Rand Miller's design principle \cite{miller2016riven}: ``The puzzle IS what the field does, not an obstacle overlaid on it.'' Scored on a five-point scale from 1 (\emph{theme is decorative; any content could substitute}) to 5 (\emph{mechanic is inseparable from the domain and could not exist in any other context}). Elegance and Reading Reward are proxies for this criterion; adding it explicitly captures residual variance not explained by the two weighted proxies. Riven Standard scores reported in Section~4.4 were inferred by the research team from existing C6 review texts rather than produced through independent blind scoring; a formal inter-rater reliability study is reserved for future work.

The six-dimension equal-weight version (/30, pass $\geq$ 22) is used in Experiments~0--2 as the controlled baseline rubric. The seven-dimension weighted version (/45, pass $\geq$ 33, with Elegance and Reading Reward each counted twice) is validated in Section~4.4 against 15 real puzzle hunt puzzles spanning three quality tiers; it achieves 24.5 percentage-point tier separation versus 10.9pp for the equal-weight baseline. Both thresholds correspond to approximately 73\% of the maximum, calibrated to the quality level at which puzzles in well-regarded MIT Mystery Hunt and Microsoft Puzzlehunt rounds have shipped. This threshold was calibrated against the authors' prior experience with puzzle hunt quality and has not been validated against a held-out set of human expert ratings. All pass/fail verdicts reported in this paper should be interpreted as valid within this author-constructed framework.

The three-reviewer design provides lens diversity (structural rigor, narrative/accessibility, domain specificity) while remaining computationally tractable. Because all reported scores represent single evaluation runs at temperature 1.0, reported scores should be understood as point estimates without sampling variance. The score ranges produced by the three-reviewer panel provide a partial proxy for evaluative uncertainty; we report these ranges when spread exceeds two points on any dimension. High variance signals genuine evaluative disagreement that averaging would suppress.

\subsection{Testbed: Puzzle Hunt Design}

A \emph{puzzle hunt} is a collection of puzzles, each with a unique verifiable answer (typically a word or short phrase), organized so that earlier answers unlock later puzzles and feeder answers combine to solve a final \emph{meta-puzzle}. Hunt puzzles are distinguished from crosswords or trivia by their use of domain knowledge as a \emph{load-bearing mechanism}: the puzzle's core step requires understanding a specific field --- a game, a musical corpus, a fictional universe, a historical event --- and cannot be completed by general reasoning alone. The solver must know the domain to solve the puzzle; the puzzle structure verifies that they do.

This makes puzzle hunt design an unusually tractable testbed for creative evaluation research, for three reasons. First, answers are objectively verifiable: there is a single correct extraction, and whether it produces the intended answer word is a binary check. This removes the most common objection to AI creative evaluation --- the absence of ground truth. Second, the expert practitioner community is documented and its evaluative language is on record: published interviews, design postmortems, and instructional texts by constructors such as Thomas Snyder, Mike Selinker, Rand Miller, Wei-Hwa Huang, Patrick Berry, and Dan Katz provide ground truth for the evaluative frameworks our profiles encode. Third, the domain admits natural variation across hunt types, allowing cross-domain generalizability studies without leaving the testbed.

We use three hunt scenarios as our experimental corpus. \textbf{Age of Empires} (``Wololo'') is a game-knowledge domain hunt built around the \emph{Age of Empires} real-time strategy series, with 5 feeder puzzles and a meta; puzzles require knowledge of civilizations, units, buildings, and game mechanics from the series. \textbf{Wavelength} is a music domain hunt with 6 puzzles requiring identification of songs, artists, albums, and musical relationships. \textbf{Dead Reckoning} is a science-fiction instrument-panel domain hunt with a projected corpus of 15 puzzles, set in a fictional starship universe with canonical reference data maintained in a locked world file; 3 puzzles were available at the time of writing with additional puzzles in progress. Puzzles, not hunts, are the unit of analysis throughout our experiments: each puzzle is evaluated independently by its three-reviewer panel, and aggregate hunt-level scores are computed from puzzle-level results.

\subsection{Experimental Design}

We report results from three experiments, one of which has complete data and two of which are in progress at the time of submission.

\textbf{Experiment~0 (Profile Ablation)} is the primary empirical contribution of this paper. We select three puzzles from the Age of Empires corpus (AoE-I, AoE-II, AoE-III, corresponding to the Dark Age, Feudal Age, and Castle Age rounds) and evaluate each under three conditions. \emph{Condition~0} is a bare prompt: the model receives the puzzle text and the instruction ``You are an expert reviewer. Score this puzzle on the following 6 dimensions.'' No persona context is provided. \emph{Condition~1} is the v1 profile panel: three reviewers are selected from the 12 v1 profiles and their abbreviated persona descriptions are loaded as system context. \emph{Condition~2} is the v2 profile panel: three reviewers are selected from the 29 v2 profiles with full 80--120-line profiles as system context. The three conditions share identical puzzle inputs, identical scoring rubrics, and identical synthesis prompts; only the profile context varies. We measure four outcomes: per-dimension score, feedback word count, named analytical frameworks introduced per review (zero in Condition~0 by construction; counted automatically from review text in Conditions~1 and~2), and actionable revision suggestions per review (manually coded into specific/general/absent). Figure~2 will present score by condition across all six dimensions; Table~3 will report full per-puzzle, per-reviewer breakdowns.

\textbf{Experiment~1 (Profile Upgrade)} is the Condition~1-to-Condition~2 arm of Experiment~0, for which complete data is already available. We report these results separately in Section~4 to allow direct comparison with the published v1 evaluation baseline, noting that Experiment~1 results are a strict subset of the Experiment~0 design. The key finding --- a +0.57 average score delta from v1 to v2 across the three puzzles, with materially larger gains in feedback specificity --- is discussed in detail in Section~4.1.

\textbf{Experiment~2 (Cross-Hunt Generalizability)} extends the v2 panel evaluation to 25 puzzles across all three hunt types: the five Age of Empires puzzles (AoE-I through AoE-V), six Wavelength puzzles, and a subset of Dead Reckoning puzzles available at evaluation time. All puzzles are evaluated under Condition~2 only. The primary questions are whether the evaluation quality gain observed in the Age of Empires corpus generalizes to music and science-fiction domain puzzles, and which named analytical frameworks transfer across domain types. Results will be reported as score deltas relative to each puzzle's v1 baseline where available, and as raw v2 scores for puzzles with no prior evaluation. Figure~4 will present score by hunt type; Table~4 will summarize hunt-level results. Data collection is ongoing; results presented in Section~4.2 reflect available data at submission time.

\textbf{Experiment~3 (Same-Input Divergence)} addresses a different question: given identical world-building inputs, does the pipeline produce structurally similar or divergent hunts across independent runs? We run the full puzzle hunt design pipeline twice from the same Age of Empires world file, producing two independent puzzle pools, round structures, and meta designs, then compare the two outputs along structural dimensions: pool selection overlap, round organization similarity, meta answer, and extraction mechanism variety. The aim is to characterize the degree of creative determinism in the pipeline --- whether panel evaluation converges on the same design choices when given the same inputs, or surfaces meaningfully different but equally valid design paths. Figure~5 will illustrate the divergence diagram; Section~4.3 will discuss implications for using AI panels as a generative design tool rather than purely an evaluative one. This experiment is preliminary; data collection is ongoing at submission time.

\textbf{Experiment~4 (Rubric Validation)} evaluates the effect of rubric design on tier separation quality, independently of profile depth. Using existing per-dimension scores from the C6 condition of Experiment~3, we recompute puzzle scores under three rubric configurations: (a) six-dimension equal-weight (/30, the baseline); (b) seven-dimension with Elegance and Reading Reward double-weighted, no Riven Standard (/40, designated C9); and (c) seven-dimension with Elegance and Reading Reward double-weighted plus Riven Standard as a scored dimension (/45, designated C11). The Riven Standard scores for (c) are inferred from the evaluative language in the C6 review texts. We measure tier separation as the percentage-point gap between the MIT/Elite tier average and the Games Magazine tier average, normalized to scale. No re-evaluation is required; this experiment operates on the existing score data from Experiment~3. Results are reported in Section~4.4.

We note two methodological limitations here, deferring full treatment to Section~5. First, our profiles encode evaluative frameworks attributed to named human experts, but the resulting AI-simulated reviews are not validated against reviews produced by those humans; we cannot directly measure the fidelity of the simulation. Second, the pass threshold of 22/30 was calibrated against the authors' prior experience with hunt quality, not against a held-out set of human expert ratings; a human expert validation study remains future work.
